\chapter{Particle Production and the Unruh Effect}
\label{ch:particle-production-unruh}
With the general formalism in place, we study two concrete examples
that isolate the essential physics of particle production in the
simplest possible settings, before tackling the full black hole
problem. The moving mirror model shows that particle production can
occur even in flat spacetime, driven purely by a time-dependent
boundary condition. The Unruh effect shows that particle production is
observer-dependent: the Minkowski vacuum contains a thermal
distribution of particles as seen by a uniformly accelerating observer.
Both of these results are more than analogies. The logarithmic relation
between the in and out coordinates in the mirror problem is precisely
the same mathematical structure that appears near a black hole horizon,
and the Unruh temperature $T = a/2\pi$ is the local limit of the
Hawking temperature experienced by a static observer near the horizon.
Understanding these simpler cases in detail makes the derivation of
Hawking radiation in the next chapter both natural and transparent.

\section{The Moving Mirror}

We now consider a concrete example of particle production in flat
spacetime, where the non-trivial physics comes not from curvature but
from boundary conditions. Consider a massless scalar field in 2D flat
spacetime with metric
\begin{equation}
	ds^2 = -dt^2 + dx^2 = -\,du\,dv,
	\qquad u = t-x,\quad v = t+x.
\end{equation}
Half the space is bounded by a mirror with trajectory $x = z(t)$. The
field satisfies a Dirichlet boundary condition at the mirror,
\begin{equation}
	\phi\big|_{x=z(t)} = 0,
\end{equation}
and the field equation $\partial^2\phi/\partial u\,\partial v = 0$
has the general solution $\phi = f(u) + g(v)$, a sum of left-moving
and right-moving waves. The boundary condition enforces a relation
between $f$ and $g$, so that the motion of the mirror mixes the two.

% [FIGURE: moving mirror with incoming and reflected modes]

We choose in-modes that in the past were standing waves,
\begin{equation}
	\phi^{\text{in}}_\omega
	= \frac{1}{\sqrt{4\pi\omega}}
	\left(e^{-i\omega v} - e^{-i\omega u}\right),
	\qquad t < 0,
\end{equation}
where $e^{-i\omega v}$ is the incoming wave and $e^{-i\omega u}$ is
the reflected wave. On the boundary $u = v$ and the two terms cancel,
satisfying the Dirichlet condition. For $t > 0$, the mirror has moved
and the reflected wave acquires a phase,
\begin{equation}
	\phi^{\text{in}}_\omega
	= \frac{1}{\sqrt{4\pi\omega}}
	\left(e^{-i\omega v} - e^{-i\omega v_0(u)}\right),
	\qquad t > 0,
\end{equation}
where $v_0(u)$ is determined by the mirror trajectory. If $\tau_u$ is
the proper time at which the mirror is at position $u = \tau_u -
	z(\tau_u)$, then $v_0 = \tau_u + z(\tau_u) = 2\tau_u - u$, giving
\begin{equation}
	\phi^{\text{in}}_\omega
	= \frac{1}{\sqrt{4\pi\omega}}
	\left(e^{-i\omega v} - e^{-i\omega(2\tau_u - u)}\right).
\end{equation}
This mode is natural for an observer in the past (simple $e^{-i\omega v}$
dependence) but complicated for an observer in the future. We therefore
also define out-modes natural for the future observer,
\begin{equation}
	\phi^{\text{out}}_{\tilde\omega}
	= \frac{1}{\sqrt{4\pi\tilde\omega}}
	\left(e^{-i\tilde\omega u}
	- e^{-i\tilde\omega(2\tau_v - v)}\right),
\end{equation}
where $\tau_v + z(\tau_v) = v$. The Bogoliubov coefficient $\beta$ is
computed on the surface $t = 0$ where the in-modes are still simple,
\begin{equation}
	\beta_{\tilde\omega\omega}
	= -\langle\phi^{\text{out}}_{\tilde\omega},
	(\phi^{\text{in}}_\omega)^*\rangle
	= -i\int_0^\infty dx\left(
	\phi^{\text{out}}_{\tilde\omega}\,
	\dot\phi^{\text{in}*}_\omega
	- \dot\phi^{\text{out}}_{\tilde\omega}\,
	\phi^{\text{in}*}_\omega\right)\bigg|_{t=0}.
\end{equation}
On this surface $u = -x$ and $v = x$, giving
\begin{equation}
	\beta_{\tilde\omega\omega}
	= -\int_0^\infty\frac{dx}{2\pi}
	\sqrt{\frac{\omega}{\tilde\omega}}
	\left(e^{i(\tilde\omega+\omega)x}
	+ e^{i(\tilde\omega-\omega)x - 2i\tilde\omega\tau_x}\right),
	\label{eq:beta-mirror}
\end{equation}
where $\tau_x + z(\tau_x) = x$. The number of produced particles is
\begin{equation}
	\frac{dN}{d\tilde\omega}
	= \int_0^\infty|\beta_{\tilde\omega\omega}|^2\,d\omega.
\end{equation}

\subsection*{Constant Velocity Mirror}

When the mirror moves with constant velocity $z = -\alpha t$,
$|\alpha| < 1$, we have $\tau_v = v/(1-\alpha)$ and so
$\tau_x = x/(1-\alpha)$. Substituting into \eqref{eq:beta-mirror} and
regularizing the oscillating integrals by $e^{-\epsilon x}$,
\begin{equation}
	\beta_{\tilde\omega\omega}
	= \frac{\sqrt{\tilde\omega\omega}}{2\pi}
	\left(\frac{2\alpha}{1-\alpha}\right)
	\frac{1}{(\tilde\omega+\omega)}
	\frac{1}{\left(\frac{1+\alpha}{1-\alpha}\tilde\omega
		+\omega\right)},
\end{equation}
where we used $\int_0^\infty dx\,e^{iAx-\epsilon x} = i/A$. Integrating
$|\beta_{\tilde\omega\omega}|^2$ over $\omega$,
\begin{equation}
	\boxed{\frac{dN}{d\tilde\omega}
		= \frac{1}{\pi^2\tilde\omega}\left(
		\frac{1}{4\alpha}\ln\frac{1+\alpha}{1-\alpha}
		- \frac{1}{2}\right).}
\end{equation}
This result is finite but non-thermal, reflecting the fact that a
uniformly moving mirror produces particles but not a thermal spectrum.

\subsection*{Accelerating Mirror}

Consider now a mirror that starts at rest, accelerates, and
asymptotically approaches a null ray, with trajectory
\begin{equation}
	z = B - t - Ae^{-2Kt}, \qquad t\to\infty.
\end{equation}
Tracing out-modes backward in time, the reflected waves never go beyond
the point $v = B$. The relation $\tau_v = -(1/2K)\ln((B-v)/A)$ for
$v < B$ gives
\begin{equation}
	\beta_{\tilde\omega\omega}
	\simeq -\frac{1}{2\pi}\sqrt{\frac{\omega}{\tilde\omega}}
	\int_0^\infty dy\,
	e^{i(\tilde\omega-\omega)(B-y)}
	\left(\frac{y}{A}\right)^{i\tilde\omega/K},
\end{equation}
where we substituted $y = B - x$ and extended the upper limit to
infinity since only the behavior near $y = 0$ matters. Using the
Gamma function integral $\Gamma(z) = \int_0^\infty dy\,y^{z-1}e^{-y}$,
\begin{equation}
	|\beta_{\tilde\omega\omega}|^2
	= \frac{1}{(K\pi)^2}\frac{\omega}{\tilde\omega}
	\,e^{-\pi\tilde\omega/2K}
	\,\Gamma\!\left(1+\frac{i\tilde\omega}{K}\right)
	\Gamma\!\left(1-\frac{i\tilde\omega}{K}\right).
\end{equation}
Using $\Gamma(z)\Gamma(1-z) = \pi/\sin\pi z$ and
$(-i)^{-i\tilde\omega/K} = e^{-\pi\tilde\omega/2K}$,
\begin{equation}
	|\beta_{\tilde\omega\omega}|^2
	= \frac{1}{2\pi K}\frac{\omega}{(\omega-\tilde\omega)^2}
	\frac{1}{e^{2\pi\tilde\omega/K}-1}.
\end{equation}
Integrating over $\omega$ and isolating the divergence at high
frequency,
\begin{equation}
	\frac{dN}{d\tilde\omega}
	= \frac{1}{2\pi K}\frac{1}{e^{2\pi\tilde\omega/K}-1}
	\int\frac{d\omega}{\omega}.
\end{equation}
The prefactor is a Bose-Einstein distribution with temperature
\begin{equation}
	T = \frac{K}{2\pi}.
\end{equation}
The integral diverges, but this represents a continuous flux rather
than a divergent total number. If the mirror stops after retarded time
$u_f$, the integral cuts off at $\omega_{\max}\sim e^{Ku_f}/A$ and
\begin{equation}
	\frac{dN}{d\tilde\omega\,du_f}
	= \frac{1}{2\pi}\cdot\frac{1}{e^{2\pi\tilde\omega/K}-1},
\end{equation}
a steady thermal flux at temperature $T = K/2\pi$. This is the
prototype of Hawking radiation, as we shall see in the next chapter.

\section{The Unruh-DeWitt Detector}

We now set up a general framework for understanding what a moving
observer actually measures. Consider a detector with internal
Hamiltonian $H_D$, following a worldline $x^\mu(\tau)$ and coupled to
a scalar field $\varphi(x)$ via the interaction
\begin{equation}
	S_{\text{int}} = \lambda\int_{-\infty}^{\infty}d\tau\,
	\varphi(x(\tau))\,\hat D(\tau),
\end{equation}
where $\hat D(\tau) = e^{iH_D\tau}\hat D(0)e^{-iH_D\tau}$ is the
detector operator in the interaction picture and $\lambda$ is the
coupling strength. To first order in perturbation theory, the
transition amplitude from the ground state $|E_0,0\rangle$ to an
excited state $|E,\psi\rangle$ is
\begin{equation}
	\mathcal{A} = i\lambda\int_{-\infty}^{\infty}d\tau\,
	\langle E|\hat D(\tau)|E_0\rangle\,
	\langle\psi|\varphi(x(\tau))|0\rangle.
\end{equation}
Summing $|\mathcal{A}|^2$ over all final field states $|\psi\rangle$
and using completeness,
\begin{equation}
	\sum_\psi\langle 0|\varphi(x(\tau_1))|\psi\rangle
	\langle\psi|\varphi(x(\tau_2))|0\rangle
	= \langle 0|\varphi(x(\tau_1))\varphi(x(\tau_2))|0\rangle
	\equiv G^+(x(\tau_1),x(\tau_2)),
\end{equation}
the total transition probability is
\begin{equation}
	P_{EE_0} = \lambda^2
	\left|\langle E|\hat D(0)|E_0\rangle\right|^2
	\int_{-\infty}^{\infty}d\tau_1\int_{-\infty}^{\infty}d\tau_2\,
	e^{-i\Delta E(\tau_1-\tau_2)}\,
	G^+(x(\tau_1),x(\tau_2)),
\end{equation}
where $\Delta E = E - E_0$. The detector response function is defined
as
\begin{equation}
	F(\Delta E) = \int_{-\infty}^{\infty}d\tau_1
	\int_{-\infty}^{\infty}d\tau_2\,
	e^{-i\Delta E(\tau_1-\tau_2)}\,
	G^+(x(\tau_1),x(\tau_2)).
\end{equation}
For stationary trajectories, $G^+$ depends only on the proper time
difference $\tau = \tau_1 - \tau_2$, and changing variables to
$\tau = \tau_1 - \tau_2$ and $T = (\tau_1+\tau_2)/2$ shows that the
transition rate per unit proper time is
\begin{equation}
	\frac{dP_{EE_0}}{d\tau}
	= \lambda^2\left|\langle E|\hat D(0)|E_0\rangle\right|^2
	\int_{-\infty}^{\infty}d\tau\,e^{-i\Delta E\tau}\,g(\tau),
\end{equation}
where $g(\tau) = G^+(x(\tau_1),x(\tau_2))$ depends only on
$\tau_1-\tau_2$. The response function $F(\Delta E)$ is therefore the
Fourier transform of the Wightman function evaluated along the
detector's worldline.

\section{The Accelerating Detector and the Unruh Effect}

We now apply the detector formalism to a uniformly accelerating
observer in 4D Minkowski spacetime. The trajectory is
\begin{equation}
	x = y = 0, \qquad z = \sqrt{t^2+B^2},
\end{equation}
with parametric form
\begin{equation}
	t = B\sinh\frac{\tau}{B}, \qquad z = B\cosh\frac{\tau}{B},
\end{equation}
and constant proper acceleration $a^\mu a_\mu = 1/B^2$.

% [FIGURE: Rindler wedge with accelerating observer trajectory]

The Wightman function in the Minkowski vacuum is
\begin{equation}
	G^+(x,x') = \langle 0|\varphi(x)\varphi(x')|0\rangle
	= \frac{1}{4\pi^2}
	\frac{1}{(t-t'-i\varepsilon)^2 - |\mathbf{x}-\mathbf{x}'|^2}.
\end{equation}
Substituting the accelerating trajectory, one finds
\begin{equation}
	G^+(x(\tau),x(\tau'))
	= -\frac{1}{16\pi^2 B^2}
	\frac{1}{\sinh^2\!\left(\dfrac{\tau-\tau'}{2B}
		- i\varepsilon\right)}.
\end{equation}
The detector response function is therefore
\begin{equation}
	F(\Delta E) = \int_{-\infty}^{\infty}
	\left(-\frac{d\tau}{16\pi^2 B^2}\right)
	\frac{e^{-i\Delta E\tau}}
	{\sinh^2\!\left(\dfrac{\tau}{2B}-i\varepsilon\right)}.
\end{equation}
We evaluate this by contour integration. The poles of
$1/\sinh^2(\tau/2B)$ occur at $\tau/2B = i\pi n$, i.e.\
$\tau = 2\pi i Bn$ for integer $n$. For $\Delta E > 0$ we close the
contour in the lower half plane, picking up poles at
$\tau = -2\pi iBn$ for $n = 1, 2, \ldots$ (the $i\varepsilon$
prescription shifts the poles away from $n=0$),
\begin{equation}
	F(\Delta E) = -\frac{1}{16\pi^2 B^2}
	\sum_{n=1}^{\infty}(-2\pi i)\,
	\underset{\tau=-2\pi iBn}{\mathrm{Res}}\,
	\frac{e^{-i\Delta E\tau}}
	{\sinh^2\!\left(\dfrac{\tau}{2B}\right)}.
\end{equation}
Computing the residues and summing the geometric series,
\begin{equation}
	F(\Delta E) = \frac{\Delta E}{2\pi}
	\sum_{n=1}^{\infty}e^{-2\pi B\Delta E\,n}
	= \boxed{\frac{\Delta E}{2\pi}
		\cdot\frac{1}{e^{2\pi B\Delta E}-1},}
\end{equation}
which is a Planck distribution with temperature
\begin{equation}
	T_U = \frac{1}{2\pi B} = \frac{a}{2\pi},
\end{equation}
where $a = 1/B$ is the proper acceleration. This is the Unruh effect:
a uniformly accelerating observer in the Minkowski vacuum measures a
thermal bath at temperature $T_U = a/2\pi$. An inertial observer in
the same state sees no particles at all; the thermal nature of the
state is entirely an effect of the acceleration.

% [FIGURE: Rindler wedge showing the causal horizon of the accelerating
%  observer]

The Unruh effect has a deep connection to the existence of a causal
horizon for the accelerating observer. The observer in the right
Rindler wedge cannot receive signals from the left wedge, and this
causal inaccessibility is responsible for the thermal character of
the state, in exact analogy with the Hawking effect for black holes,
to which we turn in the next chapter.
